% !PRINTED 741
\weberpage{741}
Accordingly $aM$, and consequently $M$ as well, is divisible by $2^{\nu'}$,
and if we set
\weq{23}{M = 2^{\nu'+\sigma}, \quad N = 2^{\nu'-\sigma-1}}
then $\sigma$ is a non-negative integral rational number.

If we substitute this in (16), it follows that:
\[ E = 2^\sigma E' T_0. \]
Now in \S.\,197, (6) we had put
\[ E = E' D \]
and there results therefore:
\[ D = 2^\sigma T_0. \]
From this there follows for the class number factor $B$, by \S.\,199, (15)
and \S.\,200, (21):
\weq{24}{B = 2^{-\sigma} \frac{T}{T_0} = 2^{-\sigma} C,}
in which $C$ is an odd integer.

Accordingly there results, by \S.\,197, (4), for the second factor
of the class number $h_1$, which, as we have seen, is an
\gesperrt{integer}:
\weq{25}{h_1 = 2^{-\sigma} h_1' C.}
If we now assume it proved that $h_1'$ is odd, then, since
$C$ is odd as well, it follows that $\sigma = 0$ and that $h_1$
is odd; and both are therewith proved by complete
induction.

With this the results derived are simplified as
follows:

3. \gesperrt{We have}
\weq{26}{M = 2^{\nu'}, \quad D = T_0, \quad B = \frac{T}{T_0};}
\gesperrt{the class number factor $B$ is equal to the ratio
of the regulator $T$ of the units $\tau$ to the regulator $T_0$
of a fundamental system of normal units, and is an
odd integer.}

If we denote the number $B$ belonging to $m = 2^\varkappa$ by
$B_\varkappa$, then the class number $h_1$ is, as follows by
complete induction,
\weq{27}{h_1 = B_4 B_5 \ldots B_\varkappa,}
and is therefore likewise an odd number. With this, then, having
regard to \S.\,198, 1., the main theorem is proved:

A. \gesperrt{The class number in the field $\Omega_m$ is odd when
$m$ is a power of 2.}
